\documentclass[11pt,a4paper]{scrartcl}

% ── Seitenlayout ──────────────────────────────────────────────────────────────
\usepackage[a4paper, top=20mm, bottom=20mm, left=22mm, right=22mm]{geometry}

% ── Typografie ────────────────────────────────────────────────────────────────
\usepackage{lmodern}
\usepackage[T1]{fontenc}
\usepackage[utf8]{inputenc}
\usepackage[ngerman]{babel}
\usepackage{microtype}
\usepackage{parskip}
\usepackage{enumitem}
\usepackage[hidelinks]{hyperref}

% ── Kopfzeile ─────────────────────────────────────────────────────────────────
\usepackage{scrlayer-scrpage}
\pagestyle{scrheadings}
\chead{}
\ohead{\small\textit{Auswertung -- Constrained Benchmarks}}
\cfoot{\pagemark}

% ═════════════════════════════════════════════════════════════════════════════
\begin{document}

\begin{center}
  {\LARGE\bfseries Auswertung der speicherbeschränkten\\RocksDB-Benchmarks}\\[6pt]
  {\large Vorläufige Analyse zur Besprechung mit dem Betreuer}\\[4pt]
  {\normalsize 4.~März~2026}
\end{center}

\bigskip

% ═════════════════════════════════════════════════════════════════════════════
\section{Was wurde gemacht}

Die zentrale Fragestellung dieser Bachelorarbeit lautet, ob RocksDB als
Out-of-Core-Zustandsspeicher für NebulaStream geeignet ist -- eine
Streamverarbeitungs-Engine für IoT- und Smart-City-Szenarien, bei denen der
verfügbare Arbeitsspeicher auf Edge-Geräten stark begrenzt ist. Um diese Frage
zu beantworten, haben wir eine Microbenchmark-Suite entwickelt, die das
RocksDB-eigene Werkzeug \texttt{db\_bench} unter kontrollierten Bedingungen
ausführt und systematisch zentrale Tuning-Parameter variiert.

Im vorliegenden ersten Experiment unter Speicherbeschränkung wurden 100
Benchmark-Konfigurationen in einem Docker-Container ausgeführt, der auf
256\,MB RAM begrenzt war (Swap deaktiviert, ein einzelner CPU-Kern). Dies
simuliert die Ressourcenausstattung eines Raspberry Pi Zero~W -- die
Geräteklasse, auf die NebulaStream für die Edge-Verarbeitung abzielt. Der
Datensatz besteht aus 2~Millionen Schlüssel-Wert-Paaren (16-Byte-Schlüssel,
100-Byte-Werte) mit einer Gesamtgröße von ca.\ 230\,MB auf der Festplatte
-- bewusst größer als der verfügbare Arbeitsspeicher, um RocksDB in den
Out-of-Core-Betrieb zu zwingen, bei dem Daten von der Festplatte gelesen und
auf diese geschrieben werden müssen, statt alles im RAM vorzuhalten.

Drei Parameter, die die Speicherverteilung innerhalb von RocksDB direkt
steuern, wurden variiert:

\begin{itemize}
  \item \textbf{write\_buffer\_size} (4, 16, 32, 64, 128\,MB) -- bestimmt, wie
        viele Daten im In-Memory-Schreibpuffer (Memtable) angesammelt werden,
        bevor sie auf die Festplatte geschrieben (geflusht) werden. Dies ist der
        größte einzelne Speicherverbraucher innerhalb von RocksDB.
  \item \textbf{cache\_size} (8, 16, 32, 64, 128\,MB) -- steuert den
        Block-Cache, der häufig gelesene Festplattenblöcke im RAM vorhält und
        damit direkt die Leseleistung beeinflusst.
  \item \textbf{compression\_type} (none, snappy, lz4, zstd) -- legt fest, ob
        und wie Daten vor dem Schreiben auf die Festplatte komprimiert werden.
        Kompression tauscht CPU-Zyklen gegen reduzierte Festplatten-I/O.
\end{itemize}

Zwei sekundäre Parameter (\texttt{level0\_slowdown\_writes\_trigger} und
\texttt{min\_level\_to\_compress}) wurden auf ihren Standardwerten fixiert, um
das Experiment auf die drei oben genannten Dimensionen zu fokussieren. Als
Workload wurde \texttt{readwhilewriting} verwendet, der gleichzeitig zufällige
Lese- und Schreibzugriffe durchführt -- dies modelliert ein
Streamverarbeitungs-Szenario, in dem die Engine gleichzeitig neue Sensordaten
aufnimmt und bestehenden Zustand nachschlägt (z.\,B.\ für Fenster-Joins oder
Aggregationen).

Von den 100 Durchläufen wurden 58 erfolgreich abgeschlossen und 42 vom
Linux-OOM-Killer beendet, als der Speicherverbrauch von RocksDB das
256\,MB-Containerlimit überschritt.


% ═════════════════════════════════════════════════════════════════════════════
\section{Was wir herausgefunden haben}

\subsection{Arbeitsspeicher ist die bindende Beschränkung, und write\_buffer\_size dominiert}

Das auffälligste Ergebnis ist das OOM-Kill-Muster. Jeder einzelne Durchlauf
mit einem 128\,MB-Schreibpuffer wurde beendet, unabhängig von Cache-Größe oder
Kompression. Bei einem 64\,MB-Schreibpuffer wurden 85\,\% der Durchläufe
abgebrochen. Der Grund ist nachvollziehbar: RocksDB hält mindestens zwei
Memtables vor (eine aktive, eine wird gerade geflusht), sodass ein
64\,MB-Schreibpuffer allein mindestens 128\,MB belegt -- die Hälfte des
gesamten Gerätespeichers -- noch bevor Block-Cache, Compaction-Puffer oder das
Betriebssystem berücksichtigt werden.

Der sichere Betriebsbereich liegt bei write\_buffer $\leq$\,32\,MB kombiniert
mit cache $\leq$\,64\,MB. Dies ist eine harte Grenze, kein gradueller Abfall:
Konfigurationen wechseln innerhalb eines einzigen Parameterschritts von
\glqq funktioniert zuverlässig\grqq{} zu \glqq stürzt immer ab\grqq.

Für NebulaStream bedeutet dies, dass der Speicherbedarf des Storage-Backends
sorgfältig gegen den Speicherbedarf der Stream-Operatoren abgewogen werden
muss. Auf einem 256\,MB-Gerät kann RocksDB realistisch höchstens 64--96\,MB
beanspruchen (Schreibpuffer + Cache zusammen), der Rest verbleibt für
NebulaStreams Operator-Zustand, Netzwerkpuffer und das Betriebssystem.

\subsection{Kompression ist nicht optional -- sie ist essenziell}

Unter unkonstrained Bedingungen (unsere früheren Baseline-Läufe auf macOS mit
unbegrenztem Speicher) machte Kompression kaum einen Unterschied, da der
gesamte Datensatz in den Arbeitsspeicher passte und Festplatten-I/O
vernachlässigbar war. Unter der 256\,MB-Beschränkung kehrt sich das Bild
vollständig um:

\begin{itemize}
  \item Unkomprimierter Speicher lieferte im Durchschnitt nur
        1.998~ops/sec, mit einer p99-Latenz von fast 12\,ms und einer
        Schreibverstärkung von 3,9$\times$. Die Engine verbrachte über 50\,\%
        ihrer Zeit im Stall-Zustand.
  \item \texttt{zstd}-Kompression lieferte 3.627~ops/sec
        (1,8$\times$ schneller), mit einer p99-Latenz unter 3,2\,ms und einer
        Schreibverstärkung von nur 1,8$\times$.
  \item \texttt{lz4} und \texttt{snappy} lagen zwischen diesen Extremen,
        wobei \texttt{lz4} \texttt{snappy} leicht übertraf.
\end{itemize}

Der Mechanismus ist klar: Kompression verkleinert die auf die Festplatte
geschriebenen und von ihr gelesenen Daten, und genau die Festplatte ist der
Engpass unter Speicherdruck. Obwohl Kompression CPU-Zyklen kostet -- und wir
nur einen einzelnen Kern haben -- überwiegen die I/O-Einsparungen den
CPU-Mehraufwand bei Weitem. Dies ist ein Schlüsselergebnis für NebulaStream:
Jeder Produktiveinsatz auf Edge-Hardware \emph{muss} Kompression im
Storage-Backend aktivieren.

Dass \texttt{zstd} (der CPU-intensivste Codec) leichtgewichtige Alternativen
wie \texttt{lz4} übertrifft, ist bemerkenswert. Es deutet darauf hin, dass
unter extremem Speicherdruck die I/O-Last so dominant ist, dass selbst
aggressive Kompression sich lohnt. Auf leistungsfähigeren Edge-Geräten mit
mehreren Kernen wäre der Vorteil von \texttt{zstd} wahrscheinlich noch größer,
da die Kompression auf einem dedizierten Kern laufen könnte.

\subsection{Ein 17-facher Leistungsabfall innerhalb derselben Hardware}

Die beste Konfiguration (write\_buffer=4\,MB, lz4, cache=8\,MB) erreichte
8.900~ops/sec. Die schlechteste überlebende Konfiguration
(write\_buffer=16\,MB, keine Kompression, cache=16\,MB) schaffte nur
518~ops/sec. Das ist ein \textbf{17-facher Unterschied} auf identischer
Hardware mit identischen Daten, verursacht ausschließlich durch die Wahl der
Konfigurationsparameter.

Dies ist wohl das wichtigste Ergebnis für die Arbeit: Es zeigt, dass die
\emph{Wahl der Tuning-Parameter wichtiger ist als die Hardware-Spezifikation}
für RocksDB auf beschränkten Geräten. Ein gut konfiguriertes 256\,MB-Gerät
übertrifft ein schlecht konfiguriertes um eine Größenordnung. Dies hat direkte
Auswirkungen auf NebulaStreams Standardkonfiguration und jede
Auto-Tuning-Strategie, die es einsetzen könnte.

\subsection{Write-Stalls sind gravierend und würden die Streamverarbeitung stören}

RocksDB blockiert eingehende Schreibvorgänge, wenn die interne Compaction nicht
mit der Ingestion-Rate Schritt halten kann. In unseren Experimenten lag der
durchschnittliche Stall-Anteil bei 34,3\,\%, mit einem Höchstwert von
71,8\,\%. Das bedeutet, dass die Storage-Engine im Durchschnitt mehr als ein
Drittel der gesamten Laufzeit neue Schreibvorgänge blockierte.

Im Kontext der Streamverarbeitung ist dies kritisch: Ein Write-Stall im
Zustandsspeicher propagiert als Gegendruck (Back-Pressure) durch die gesamte
Operator-Pipeline. Wenn die Storage-Engine für 100\,ms blockiert, kann kein
vorgelagerter Operator in dieser Zeit Ergebnisse ausgeben. Für
Echtzeit-Anwendungen (z.\,B.\ Ampelsteuerung auf Basis von
Verkehrskameradaten) kann das verpasste Deadlines bedeuten.

Wir beobachteten ein kontraintuitives Muster: Der kleinste Schreibpuffer
(4\,MB) hatte den \emph{höchsten} Stall-Anteil (40--63\,\%), aber auch den
\emph{höchsten} Durchsatz. Die Erklärung: Kleine Puffer flushen häufig (viele
kurze Stalls), während große Puffer selten flushen, aber jeder einzelne Flush
unter Speicherdruck teuer ist. Der Nettoeffekt ist, dass viele kleine Flushes
schneller abgeschlossen werden als wenige große, wenn der Speicher knapp ist.

\subsection{Cache-Trefferquote skaliert linear, aber mit abnehmendem Grenznutzen}

Die Block-Cache-Trefferquoten reichten von 4,5\,\% bei 8\,MB bis 64,1\,\% bei
128\,MB und skalierten annähernd linear mit der Cache-Größe. Es gab keinen
plötzlichen Abfall oder Wendepunkt. Ein 32\,MB-Cache (12,5\,\% des
Gerätespeichers) erreichte eine Trefferquote von 18,8\,\%, was wir als
vernünftigen Ausgangspunkt betrachten, der ausreichend Speicher für den
Schreibpuffer und NebulaStreams eigene Operatoren übrig lässt.


% ═════════════════════════════════════════════════════════════════════════════
\section{Einordnung in die Bachelorarbeit}

\subsection{Umfang und Einschränkungen dieser Ergebnisse}

Es ist wichtig, transparent zu benennen, was dieses Experiment aussagen kann
und was nicht. Wir haben einen \textbf{schmalen Ausschnitt} des gesamten
Parameterraums getestet:

\begin{itemize}
  \item \textbf{Ein einzelner Workload} (\texttt{readwhilewriting}) von
        mehreren relevanten. Reine Schreiblast (\texttt{fillrandom}), reine
        Leselast (\texttt{readrandom}) und Zustandsbereinigung
        (\texttt{deleterandom}) wurden noch nicht getestet.
  \item \textbf{Drei von fünf Tuning-Dimensionen.} Der L0-Compaction-Trigger
        und das minimale Kompressionslevel wurden fixiert. Diese Parameter
        könnten mit Schreibpuffer- und Cache-Größe in Weisen interagieren, die
        wir noch nicht erfasst haben.
  \item \textbf{Synthetische Schlüssel-Wert-Größen.} Unsere 16-Byte-Schlüssel
        und 100-Byte-Werte sind eine vertretbare Annäherung, entsprechen aber
        nicht den tatsächlichen NebulaStream-Zustandsformaten, die je nach
        Operatortyp variieren.
  \item \textbf{Simulierte Hardware.} Dockers cgroup-Speicherlimits auf macOS
        erzwingen die 256\,MB-Grenze zuverlässig, aber der zugrunde liegende
        Speicher (SSD über die Docker-VM) ist schneller als die SD-Karten
        realer Raspberry-Pi-Geräte. Unsere Durchsatzzahlen sind daher
        \emph{optimistische Obergrenzen} -- tatsächliche Edge-Hardware wäre
        langsamer.
  \item \textbf{Keine Konkurrenz durch NebulaStream selbst.} In der Produktion
        konkurriert das Storage-Backend um CPU und Speicher mit NebulaStreams
        Query-Engine, Netzwerk-Stack und weiteren Operatoren. Unser Benchmark
        betreibt RocksDB isoliert, sodass der Speicher- und CPU-Spielraum
        großzügiger ist als in der Realität.
\end{itemize}

Aufgrund dieser Einschränkungen sollten die Ergebnisse als
\textbf{Richtungsbefunde} interpretiert werden -- sie identifizieren, welche
Parameter wichtig sind und wo die kritischen Schwellenwerte liegen, aber die
exakten Zahlen (z.\,B.\ \glqq 8.900~ops/sec mit 4\,MB-Schreibpuffer\grqq)
lassen sich nicht direkt auf einen realen NebulaStream-Einsatz übertragen. Die
\emph{relativen} Zusammenhänge (Kompression hilft, große Schreibpuffer führen
zum Absturz, Stalls sind gravierend) sind robust und sollten über
Konfigurationen hinweg Bestand haben.

\subsection{Was wir belastbar aussagen können}

Trotz des begrenzten Umfangs sind mehrere Schlussfolgerungen gut gestützt:

\begin{enumerate}
  \item \textbf{RocksDB kann innerhalb von 256\,MB betrieben werden} -- aber
        nur, wenn Schreibpuffer und Cache konservativ dimensioniert werden
        ($\leq$\,32\,MB bzw.\ $\leq$\,64\,MB). Die Machbarkeitsfrage wird mit
        einem bedingten \glqq Ja\grqq{} beantwortet.
  \item \textbf{Kompression ist zwingend erforderlich.} Die
        Leistungseinbußen bei unkomprimiertem Speicher unter Speicherdruck
        sind so gravierend, dass Kompression eine nicht verhandelbare
        Standardeinstellung für jeden Edge-Einsatz sein muss.
  \item \textbf{Die Konfigurationssensitivität ist extrem.} Ein 17-facher
        Leistungsunterschied innerhalb desselben Speicherbudgets bedeutet, dass
        NebulaStream nicht mit naiven Standardwerten ausgeliefert werden kann --
        es braucht entweder sorgfältig gewähltes statisches Tuning oder
        adaptives Auto-Tuning.
  \item \textbf{Write-Stalls sind ein reales Problem für die
        Streamverarbeitung.} Das Storage-Backend kann die
        Verarbeitungs-Pipeline für signifikante Zeitanteile blockieren, was in
        NebulaStreams Back-Pressure- und Scheduling-Design berücksichtigt werden
        muss.
\end{enumerate}

\subsection{Was Annahme oder Schlussfolgerung bleibt}

Die folgenden Punkte sind auf Basis unserer Daten plausibel, aber noch nicht
bewiesen:

\begin{itemize}
  \item Dass \texttt{zstd} auch auf tatsächlicher Single-Core-ARM-Hardware
        optimal bleibt (unsere Benchmarks liefen auf einem x86-Docker-Host, auf
        dem \texttt{zstd}s SIMD-Optimierungen sich anders verhalten könnten als
        auf ARM).
  \item Dass sich die OOM-Grenzen vorhersagbar verschieben, wenn NebulaStreams
        eigener Speicherbedarf hinzukommt (wir haben beides nicht gemeinsam
        betrieben).
  \item Dass die Write-Stall-Muster auch für reine Ingestion
        (\texttt{fillrandom}) oder Zustandsbereinigung (\texttt{deleterandom})
        gelten, die andere I/O-Profile aufweisen.
  \item Dass Schlüssel-/Wertgrößen, die realen NebulaStream-Operatoren
        entsprechen, qualitativ ähnliche Ergebnisse liefern (größere Werte
        würden das Kompressionsverhältnis und die Cache-Effizienz verändern).
\end{itemize}


% ═════════════════════════════════════════════════════════════════════════════
\section{Vorgeschlagene nächste Schritte}

Basierend auf diesen vorläufigen Ergebnissen schlagen wir folgenden Plan für
die verbleibende Arbeit vor, grob nach Priorität geordnet:

\begin{enumerate}

  \item \textbf{\texttt{fillrandom}-Workload ergänzen} (hohe Priorität,
        geringer Aufwand).\\
        Reine sequentielle Ingestion ist das häufigste Zugriffsmuster für
        Sensordatenströme. Die Durchführung des gleichen Parameter-Sweeps mit
        \texttt{fillrandom} würde bestätigen, ob die Erkenntnisse zu
        Kompression und Schreibpuffer auf reine Schreiblasten übertragbar sind,
        und Write-Stall-Daten ohne den Störfaktor gleichzeitiger Lesezugriffe
        liefern.

  \item \textbf{Vollständigen Parameter-Sweep durchführen} (hohe Priorität,
        hoher Aufwand).\\
        Das aktuelle Experiment mit 100 Durchläufen hat zwei Parameter fixiert.
        Der vollständige Sweep (900 Kombinationen) ergänzt
        \texttt{level0\_slowdown\_writes\_trigger} (10, 20, 40) und
        \texttt{min\_level\_to\_compress} ($-1$, 0, 2). Der L0-Trigger steuert
        direkt, ab wann RocksDB Schreibvorgänge drosselt, was unter
        Speicherdruck signifikant mit der Schreibpuffergröße interagieren
        könnte. Geschätzte Laufzeit: 2--5 Tage mit parallelen Docker-Workern.

  \item \textbf{NebulaStream-realistische Schlüssel-/Wertgrößen verwenden}
        (mittlere Priorität).\\
        Die NebulaStream-Codebasis konsultieren, um tatsächliche
        Zustandsrecord-Größen gängiger Operatoren zu ermitteln (z.\,B.\
        Fenster-Count: kleiner Wert; Fenster-Join: großer Wert mit
        serialisierten Tupeln). Eine fokussierte Teilmenge der Konfigurationen
        mit diesen Größen erneut durchführen, um die externe Übertragbarkeit zu
        validieren.

  \item \textbf{\texttt{deleterandom}-Workload ergänzen} (mittlere Priorität,
        geringer Aufwand).\\
        Streamverarbeitungs-Engines lassen alten Zustand verfallen (z.\,B.\
        geschlossene Fenster, eviktierte Join-Einträge). Löschmuster
        beeinflussen das Compaction-Verhalten und die Space-Amplification anders
        als Einfügungen. Dies würde die Workload-Abdeckung abrunden.

  \item \textbf{Test auf realer Raspberry-Pi-Hardware} (niedrigere Priorität,
        hoher Aufwand).\\
        Ein kleiner Validierungslauf (10--20 Konfigurationen) auf einem
        physischen Raspberry Pi Zero~W mit SD-Karte würde die Docker-Simulation
        empirisch absichern. Das Hauptrisiko ist, dass die I/O-Latenz der
        SD-Karte Write-Stalls über das hinaus verstärkt, was wir beobachtet
        haben. Selbst ein begrenzter Lauf auf realer Hardware würde die
        Schlussfolgerungen der Arbeit erheblich stärken.

  \item \textbf{Integration mit NebulaStream-Prototyp} (Stretch-Ziel).\\
        Die ultimative Validierung wäre die Ausführung einer einfachen
        NebulaStream-Pipeline (z.\,B.\ eine gleitende Fensteraggregation über
        synthetische Sensordaten) mit RocksDB als Zustandsspeicher auf
        beschränkter Hardware. Dies geht über den Scope des Microbenchmarks
        hinaus, könnte aber eine starke Demonstration im Evaluationskapitel
        der Arbeit sein.

\end{enumerate}

\bigskip
\noindent\rule{\textwidth}{0.3pt}
\smallskip
{\small Dieses Dokument begleitet den vorläufigen Ergebnisbericht
(\texttt{report\_preliminary.tex}) und dient als Vorbereitungsmaterial für die
Betreuerbesprechung am 4.~März~2026. Der Ergebnisbericht enthält die Rohdaten,
Visualisierungen und tabellierten Befunde; dieses Dokument liefert den
interpretativen Kontext.}

\end{document}
